\section{Conclusion}\label{sec:conclusion}
This paper has proposed an energy-frugal, cross-layer task-oriented semantic communication framework for edge Visual Question Answering (VQA) over realistic wireless channels. By integrating deep variational neural source compression with standardized digital 3GPP $\text{LDPC}(1020, 512)$ channel coding and an adapted foundational Vision--Language Model (Qwen3-VL-4B), our framework bridges the critical gap between theoretical joint source-channel coding and practical digital wireless standards.

Under a rigorous green dual-mode energy formulation, we demonstrated that in energy-constrained edge uplinks (Mode~2), dynamically selecting compact source bitrates unlocks up to $+3.01\text{ dB}$ and $+6.00\text{ dB}$ of physical SNR gain, effectively preventing catastrophic packet erasure in deep fading regimes. We developed a lightweight, channel-aware semantic policy network that dynamically co-optimizes the transmitter's source bitrate and the receiver's visual token budget based on multimodal query semantics and instantaneous channel CSI. Large-scale evaluation over 2,400 independent blind test images across 10 channel SNR points on the TDIUC benchmark confirms that the proposed policy achieves a $+29.92\%$ accuracy advantage over fixed-rate baselines and a $+27.71\%$ advantage over channel-blind routing under harsh fading ($2.5\text{ dB}$). Furthermore, our framework establishes non-dominated Pareto efficiency frontiers across RF transmission energy, end-to-end latency, and task accuracy, while cutting visual token overhead by $49.2\%$ for low-entropy semantic tasks. These results underscore cross-layer semantic-channel co-design as an indispensable principle for green, resilient 6G edge intelligence.
